\subsection{Explicit nonzero projective classes for the two remaining finite groups}\label{D:explicit-projectives}
\begin{proposition}\label{D:Q8-projector}
Let $R=\Z Q_8$ and $N=\sum_{g\in Q_8}g$. The matrix
\[
 e_8=\begin{pmatrix}1-8N&21N\\-3N&8N\end{pmatrix}
\]
is idempotent, and its image has nonzero class in $\widetilde K_0(R)$.
\end{proposition}
\begin{proof}
The identity $N^2=8N$ verifies all four entries of $e_8^2=e_8$.
The row $(3,-8)$ maps its two columns to $3$ and $-N$, so it maps
$\im e_8$ onto the ideal $(3,N)\subset R$. Both groups have integral
rank eight: on every nontrivial rational representation $e_8$ is
$\operatorname{diag}(1,0)$, and on the trivial representation it is a
rank-one projector. The surjection is therefore an isomorphism.

Here is the precise external nonvanishing input and its applicability.
Martinet's Theorem I.1 \cite[pp.~400--401]{martinet1971} defines the
rank-one projective modules $P_a$ for odd $a$, proves that $P_1=R$,
and proves $P_a\cong P_b$ precisely when $a\equiv\pm b\pmod8$.
His remark on p.~401 identifies $(3,N)$ with $P_3$.
Thus $(3,N)$ is not free. His Theorem II.1, proved on pp.~402--404,
states that every finitely generated stably free $\Z Q_8$-module is
free. These hypotheses apply to the image of our explicit finite
idempotent, so its reduced class cannot vanish. This uses stable
nonfreeness, not merely nonfreeness. Ranicki's explicit projection
calculation \cite[pp.~4--5]{ranickiExplicit} gives the same matrix and
identification independently. Martinet's 2014 correction
\cite{martinet2014correction} concerns the arithmetic sign choice in
Theorem IV.1; it does not alter Theorems I.1 or II.1. We opened the
original proof and the correction.
\end{proof}

\begin{proposition}\label{D:C23-projector}
Let $R=\Z C_{23}=\Z[g]/(g^{23}-1)$, set $N=\sum_{a=0}^{22}g^a$, and put
\[
 \omega=1+g+g^2+g^3+g^4+g^6+g^8+g^9+g^{12}+g^{13}+g^{16}+g^{18}.
\]
Then
\[
 e_{23}=\begin{pmatrix}
 4+8N&\omega-1+N\\
 2\omega-46N&-3-8N
 \end{pmatrix}
\]
is idempotent and $[\im e_{23}]\ne0$ in $\widetilde K_0(R)$.
\end{proposition}
\begin{proof}
Counting sums of the displayed quadratic residues modulo $23$ gives
\[
 N^2=23N,\qquad \omega N=12N,\qquad
 \omega^2-\omega+6=6N.
\]
Substitution proves the idempotent identity entry by entry. Alternatively,
the injective ring map
\[
 R\longrightarrow \Z[\zeta_{23}]\times\Z,
 \qquad g\longmapsto(\zeta_{23},1)
\]
identifies $R$ with the pairs whose residues agree in $\F_{23}$.
The two images of $e_{23}$ are
\[
 \begin{pmatrix}4&\omega-1\\2\omega&-3\end{pmatrix},
 \qquad
 \begin{pmatrix}188&34\\-1034&-187\end{pmatrix}
   =\binom{2}{-11}(94,17).
\]
The second is idempotent since $188-187=1$. In the first ring
$\omega^2-\omega+6=0$; writing
$p=(2,\omega)^{\mathsf T}$ and $q=(2,(\omega-1)/2)$ gives $qp=1$,
so the first matrix is $pq$.

Let $L=\Q(\zeta_{23})$, $K=\Q(\omega)=\Q(\sqrt{-23})$, and
$\mathcal O_K=\Z[\omega]$. The polynomial discriminant is $-23$,
so this is the full ring of integers. The ideal
$I=(2,\omega)\subset\mathcal O_K$ has norm $2$.
It is not principal: a generator would have norm $2$, whereas
\[
 \operatorname{Norm}_{K/\Q}(a+b\omega)
 =a^2+ab+6b^2=(a+b/2)^2+23b^2/4
\]
never equals $2$ for integers $a,b$. The conjugate prime
$(2,\omega-1)$ does not contain $\omega-2$, and this element has
norm $8$. Unique ideal factorization therefore gives
$I^3=(\omega-2)$, so $[I]$ has exact order three.

The extension degree $[L:K]=11$ is prime to three. If
$I\mathcal O_L$ were principal, its ideal norm down to $K$ would
make $I^{11}$ principal, a contradiction. The image of $pq$ is
isomorphic to $(I\mathcal O_L)^{-1}$: the injection
$r\mapsto(2r,\omega r)$ has precisely that fractional ideal as its
integral domain, and $q$ is its inverse on the image because $qp=1$.
A stably free rank-one module over the commutative ring
$\mathcal O_L$ is free by taking its top exterior power. Hence the
extension of $\im e_{23}$ to $\mathcal O_L$ is not stably free,
and neither is $\im e_{23}$ over $R$.
\end{proof}

\begin{remark}\label{D:projector-check}
The program \texttt{d\_explicit\_projectives.py} gives both complete group
tables, verifies all $512$ and $12\,167$ associativity triples, checks
all integral coefficients of both projector identities, and computes
the exact ranks $8$ and $23$ of their integral regular matrices.
It also checks both complementary projectors. Nontriviality of a
projective class is established by the proofs above, not inferred from
a matrix rank calculation. These are algebraic inputs to the finite
category construction; by themselves they are not counterexamples to
Question D1.
\end{remark}
